\chapter{Chirurgia geometryczna i odwzorowania normalne}
\label{ch:chirurgia-geometryczna-21}
\index{chirurgia}

\noindent\textbf{Pytanie przewodnie.}
Rozdział~\ref{ch:rachunek-uchwytow} opisał, jak dołączenie uchwytu
zmienia brzeg kobordyzmu, a Rozdział~\ref{ch:s-kobordyzm}
wyjaśnił, kiedy pewien $h$-kobordyzm jest iloczynem.
Teraz zmieniamy samą rozmaitość w sposób kontrolowany:
usuwamy rurkowate otoczenie sfery i wklejamy drugi kawałek
o tym samym brzegu. Taką zmianę nazywamy \emph{chirurgią}.
Jej zadaniem jest modyfikowanie odwzorowania $M\to X$,
aż stanie się równoważnością homotopijną.
Operację geometryczną przeprowadzimy i policzymy szczegółowo;
przeszkody leżące w grupach $L$ należą do następnego rozdziału.

\section{Operacja: sfera, obramowanie i wymiana kawałka}
\label{sec:krok-chirurgii-21}

Niech $M^n$ będzie gładką rozmaitością zamkniętą.
Dla $0\leq p\leq n-2$ połóżmy $q=n-p\geq2$.
Wybieramy gładkie zanurzenie
\[
 e:S^p\times D^q\longrightarrow M.
\]
Jego rdzeń $e(S^p\times\{0\})$ jest sferą $p$-wymiarową.
Parametryzacja dysków $\{x\}\times D^q$ to
\emph{obramowanie} jej wiązki normalnej: wskazuje $q$
niezależnych kierunków normalnych, zmieniających się
gładko z $x$. Samo zanurzenie $S^p\hookrightarrow M$
bez takiego wyboru nie określa jeszcze chirurgii.

\begin{definicja}[Chirurgia typu $(p,q)$]
\label{def:chirurgia-21}
\emph{Wynikiem chirurgii} względem $e$ jest
rozmaitość, po wygładzeniu naroży,
\begin{equation}\label{eq:chirurgia-21}
 M_e=\bigl(M\setminus\operatorname{int}e(S^p\times D^q)\bigr)
 \mathbin{\cup_{\,S^p\times S^{q-1}}}
 (D^{p+1}\times S^{q-1}).
\end{equation}
Identyfikację na brzegu wyznacza $e$ i naturalna
identyfikacja $\partial(D^{p+1}\times S^{q-1})
=S^p\times S^{q-1}$.
Mówimy też o \emph{chirurgii na obramowanej sferze $S^p$}.
\end{definicja}

Wycinany kawałek ma brzeg $S^p\times S^{q-1}$.
Wklejany kawałek ma \emph{dokładnie ten sam brzeg},
lecz zamienia role obu kierunków: dyski normalne starej
sfery zastępuje dysk $D^{p+1}$ wypełniający tę sferę.
Wzór~\eqref{eq:chirurgia-21} jest wzorem na
brzeg dołączenia $(p+1)$-uchwytu z
Twierdzenia~\ref{thm:slad-chirurgii-18}; nie
jest to nowa, niezależna konstrukcja.

\begin{figure}[htbp]
\centering
\begin{tikzpicture}[>=Latex,font=\small,line cap=round]
 \node[align=center,manifoldblue!75!black] at (-2.55,1.44)
  {$S^0\times D^2$\\dwa dyski};
 \fill[manifoldblue!14] (-3.30,0) circle (.66);
 \fill[manifoldblue!14] (-1.80,0) circle (.66);
 \draw[manifoldblue!80!black,very thick]
  (-3.30,0) circle (.66) (-1.80,0) circle (.66);
 \node at (-2.55,-1.15) {wycinamy};
 \draw[->,very thick,manifoldred!80!black] (-.95,0)--(.44,0);
 \node[align=center,manifoldred!80!black] at (-.23,.53)
   {wspólny brzeg\\$S^0\times S^1$};
 \begin{scope}[shift={(2.32,0)}]
  \fill[even odd rule,manifoldgold!27]
      (0,0) circle (.84) (0,0) circle (.40);
  \draw[manifoldgold!75!black,very thick]
      (0,0) circle (.84) (0,0) circle (.40);
 \end{scope}
 \node[align=center,manifoldgold!65!black] at (2.32,1.44)
  {$D^1\times S^1$\\pierścień};
 \node at (2.32,-1.15) {wklejamy};
\end{tikzpicture}
\caption{Lokalny model dla $p=0$, $n=2$.
Dwa dyski i pierścień mają ten sam brzeg: dwa okręgi.
Ostateczną rozmaitość otrzymujemy, przyklejając
pierścień do części starej powierzchni pozostałej
po wycięciu dysków. Globalny przykład $S^2\to T^2$
widzieliśmy na Rysunku~\ref{fig:chirurgia-sfera-torus-18}.}
\label{fig:krok-chirurgii-21}
\end{figure}

\begin{przyklad}[Zawęźlony torus jako wynik chirurgii]
\label{prz:knotted-torus-surgery-21}
W poprzednim rysunku liczyliśmy zmianę
\emph{abstrakcyjnej} powierzchni. Możemy
również kontrolować jej zanurzenie w $\R^3$.
Niech $K\subset\R^3$ będzie gładkim węzłem
trójlistnym, na przykład obrazem krzywej
\[
 t\longmapsto(\sin t+2\sin2t,
 \cos t-2\cos2t,\sin3t),
 \qquad 0\leq t\leq2\pi.
\]
Wybierzmy małą kulę $B^3$, która zawiera
tylko krótki łuk $K$. Pozostałą część $K$
pogrubiamy do wąskiej rurki $D^1\times D^2$
i jej oba końce przyklejamy do dwóch dysków
w $\partial B^3$. Otrzymujemy pełny torus
$B^3\cup(D^1\times D^2)$ będący otoczeniem
rurkowym $N(K)$. Na jego brzegu operacja jest
dokładnie chirurgią $S^2\to T^2$ z
Rysunku~\ref{fig:krok-chirurgii-21}:
usuwamy dwa dyski $S^0\times D^2$
i wklejamy pobocznicę rurki $D^1\times S^1$.
Jeśli rdzeń uchwytu jest zawęźlony,
otrzymujemy \emph{zawęźlony torus}
$T_K=\partial N(K)\subset\R^3$.
Jako samodzielna powierzchnia nadal jest
homeomorficzny ze zwykłym $T^2$; różni się
położeniem w przestrzeni i dopełnieniem.
\end{przyklad}

\begin{figure}[htbp]
\centering
\begin{tikzpicture}[line cap=round,font=\small]
 \def\knotpath#1#2{plot[domain=#1:#2,samples=180]
  ({.48*(sin(\x)+2*sin(2*\x))},
   {.48*(cos(\x)-2*cos(2*\x))})}
 \begin{scope}[shift={(-4.15,0)}]
  \node[manifoldblue!75!black] at (0,2.05)
    {$\partial B^3=S^2$};
  \shade[inner color=white,outer color=manifoldblue!38]
    (0,0) circle (.94);
  \draw[manifoldblue!76!black,thick] (0,0) circle (.94);
  \fill[manifoldgold!43] (-.38,-.09) ellipse (.17 and .28);
  \fill[manifoldgold!43] (.38,-.09) ellipse (.17 and .28);
  \draw[manifoldgold!75!black,thick]
    (-.38,-.09) ellipse (.17 and .28)
    (.38,-.09) ellipse (.17 and .28);
  \node[font=\scriptsize] at (0,-1.42)
    {dwa dyski $S^0\times D^2$};
 \end{scope}
 \draw[-{Latex[length=2.5mm]},manifoldred!80!black,
    line width=1.1pt] (-2.90,0)--(-1.70,0);
 \begin{scope}
  \node[manifoldblue!75!black] at (0,2.05)
    {rdzeń: węzeł $K$};
  \draw[manifoldblue!76!black,line width=2.0pt]
    \knotpath{0}{360};
  \foreach \a/\b in {8/23,128/143,248/263}{
   \draw[white,line width=6.5pt] \knotpath{\a}{\b};
   \draw[manifoldblue!76!black,line width=2.0pt]
     \knotpath{\a}{\b};
  }
  \draw[manifoldgold!78!black,dashed,thick]
    (0,-.48) circle (.29);
  \node[manifoldgold!60!black,font=\scriptsize,
    anchor=north west] at (.30,-.62) {$B^3$};
 \end{scope}
 \draw[-{Latex[length=3mm]},manifoldred!80!black,
    line width=1.2pt] (1.70,0)--(3.10,0);
 \node[manifoldred!80!black,font=\scriptsize,
   align=center] at (2.40,1.30) {pogrubiamy\\łuk};
 \begin{scope}[shift={(4.55,0)}]
  \node[manifoldblue!75!black] at (0,2.05)
    {$T_K=\partial N(K)$};
  \draw[manifoldblue!75!black,line width=17pt]
    \knotpath{0}{360};
  \draw[manifoldblue!25,line width=14pt]
    \knotpath{0}{360};
  \foreach \a/\b in {8/23,128/143,248/263}{
   \draw[white,line width=21pt] \knotpath{\a}{\b};
   \draw[manifoldblue!75!black,line width=17pt]
     \knotpath{\a}{\b};
   \draw[manifoldblue!25,line width=14pt]
     \knotpath{\a}{\b};
  }
 \end{scope}
\end{tikzpicture}
\caption{Przestrzenny przykład operacji z
Rysunku~\ref{fig:krok-chirurgii-21}.
Najpierw zaznaczamy dwa dyski na sferze.
W środku widać rzut węzła trójlistnego;
białe przerwy oznaczają przejścia pod spodem,
a złoty okrąg małą kulę wokół fragmentu rdzenia.
Po prawej pogrubiona pozostała część węzła
jest $1$-uchwytem przyklejonym do kuli;
jej brzeg to zawęźlony torus.}
\label{fig:knotted-torus-surgery-21}
\end{figure}

\begin{uwaga}[Dlaczego nie powstaje rogata sfera?]
Rogata sfera Alexandra jest topologicznie
$S^2$, lecz jej zanurzenie w $\R^3$ nie jest
lokalnie płaskie na zbiorze punktów granicznych
rogów. Każdy \emph{skończony} ciąg opisanych
tutaj gładkich dołączeń uchwytów ma gładki,
a więc lokalnie płaski brzeg. Rogata sfera
powstaje jako granica nieskończonego
powtarzania zazębiających się rogów, nie
jako jeden krok gładkiej chirurgii. Dla
jej konstrukcji i analizy dopełnienia zob.
\href{https://people.mpim-bonn.mpg.de/aruray/documents/top_manifolds.pdf}
{\emph{Topological Manifolds}, §5.2}.
\end{uwaga}

\begin{uwaga}[Dlaczego obramowanie ma znaczenie?]
Z trywialności wiązki normalnej wynika \emph{możliwość}
wyboru $e$, ale nie jego jednoznaczność.
Dwa obramowania różnią się odwzorowaniem
$S^p\to O(q)$ (po ustaleniu zanurzenia rdzenia).
Wymiana kawałków w~\eqref{eq:chirurgia-21} może więc
dać różne wyniki. W szczególności dla okręgu
w $3$-rozmaitości obramowania określają liczby
całkowite, gdyż $\pi_1(SO(2))\cong\Z$;
to źródło różnych wyników chirurgii na węźle.
Jeśli zaś wiązka normalna jest nietrywialna,
nie istnieje nawet zanurzenie $S^p\times D^q$
odpowiadające jej całej tubie.
Na przykład diagonalna sfera $\Delta\subset S^2\times S^2$
ma wiązkę normalną izomorficzną z $TS^2$,
o liczbie Eulera $2$; nie można wykonać na niej
chirurgii wymagającej iloczynowego otoczenia
$S^2\times D^2$.
\end{uwaga}

\section{Ślad chirurgii, odwrotność i dualny uchwyt}
\label{sec:slad-chirurgii-21}

\begin{definicja}[Ślad chirurgii]
\index{slad chirurgii@Ślad chirurgii}
\label{def:slad-chirurgii-21}
\emph{Śladem} operacji $e$ nazywamy $(n+1)$-wymiarowy
kobordyzm
\begin{equation}\label{eq:slad-chirurgii-21}
 W_e=(M\times[0,1])\cup_{e\times\{1\}}
        (D^{p+1}\times D^q).
\end{equation}
Po wygładzeniu naroży ma on brzeg wejściowy $M$
i brzeg wyjściowy $M_e$.
\end{definicja}

\begin{twierdzenie}[Ślad i chirurgia odwrotna]
\label{thm:odwrotna-chirurgia-21}
Ślad~\eqref{eq:slad-chirurgii-21} powstaje przez
dołączenie jednego uchwytu indeksu $p+1$.
Oglądany od strony $M_e$ jest śladem chirurgii
na obramowanej sferze $S^{q-1}\subset M_e$,
o uchwycie indeksu $q=n-p$.
Ta druga chirurgia przywraca $M$.
\end{twierdzenie}
\begin{proof}
Uchwyt w~\eqref{eq:slad-chirurgii-21} ma postać
$D^{p+1}\times D^q$. Jego część przyklejona do
$M\times\{1\}$ to $\partial D^{p+1}\times D^q
=S^p\times D^q$, więc indeks uchwytu wynosi $p+1$.
Niewykorzystana część jego brzegu to
$D^{p+1}\times\partial D^q=D^{p+1}\times S^{q-1}$;
wraz z pozostałą częścią $M$ daje dokładnie
wzór~\eqref{eq:chirurgia-21}.

Po odwróceniu kierunku kobordyzmu część
$D^{p+1}\times S^{q-1}$ jest obszarem
przyczepienia \emph{dualnego} uchwytu
$D^q\times D^{p+1}$. Sferą rdzenia w nowym
brzegu jest $\{0\}\times S^{q-1}$, a obramowaniem
jej normalnych dysków jest czynnik $D^{p+1}$.
Wycinamy więc $D^{p+1}\times S^{q-1}$
i przywracamy $S^p\times D^q$.
Wymiar uchwytu dualnego wynosi
$(n+1)-(p+1)=q$, zgodnie z
Twierdzeniem~\ref{thm:dualny-uchwyt-18}.
\end{proof}

\begin{figure}[htbp]
\centering
\begin{tikzpicture}[>=Latex,font=\small,line cap=round]
 \node[draw=manifoldblue!70!black,rounded corners=3pt,
  fill=manifoldblue!8,minimum width=1.8cm,minimum height=.8cm]
  (M) at (-5.05,0) {$M^n$};
 \node[draw=manifoldgold!80!black,rounded corners=3pt,
  fill=manifoldgold!12,minimum width=2.7cm,minimum height=1.1cm,
  align=center] (W) at (0,0) {$W_e^{n+1}$\\jeden uchwyt};
 \node[draw=manifoldgreen!70!black,rounded corners=3pt,
  fill=manifoldgreen!9,minimum width=1.8cm,minimum height=.8cm]
  (N) at (5.05,0) {$M_e^n$};
 \draw[->,thick,manifoldred!75!black]
  (M) -- node[above,font=\scriptsize] {indeks $p+1$} (W);
 \draw[->,thick,manifoldred!75!black]
  (N) -- node[above,font=\scriptsize] {indeks $q=n-p$} (W);
 \node[align=center,text width=10cm] at (0,-1.5)
  {Ten sam kobordyzm ma dwa opisy uchwytowe,
   zależne od tego, który brzeg uznamy za wejściowy.};
\end{tikzpicture}
\caption{Ślad chirurgii od strony starego i nowego brzegu.
Różnica indeksów nie oznacza dwóch różnych kobordyzmów.}
\label{fig:slad-chirurgii-21}
\end{figure}

\section{Jak chirurgia zmienia niezmienniki}
\label{sec:niezmienniki-chirurgia-21}

W tej sekcji $M$ jest spójna i zamknięta.
Ślad pozwala obliczyć zmianę bez zgadywania
topologii $M_e$ z samego rysunku.

\begin{twierdzenie}[Względna homologia śladu]
\label{thm:homologia-sladu-21}
Dla śladu $W_e$ mamy
\[
 H_i(W_e,M;\Z)\cong
 \begin{cases}\Z,&i=p+1,\\0,&i\ne p+1,
 \end{cases}
 \qquad
 H_i(W_e,M_e;\Z)\cong
 \begin{cases}\Z,&i=q,\\0,&i\ne q.
 \end{cases}
\]
Odpowiadające ciągi dokładne określają wpływ
chirurgii na homologię obu brzegów.
\end{twierdzenie}
\begin{proof}
Obcięcie $M\times[0,1]$ w kierunku przedziału
nie zmienia homologii. Z twierdzenia o wycinaniu
względny kompleks pary $(W_e,M)$ ma homologię
pary
\[
 (D^{p+1}\times D^q,S^p\times D^q).
\]
Projekcja w pierwszy czynnik jest homotopijną
równoważnością par z $(D^{p+1},S^p)$:
odkształcamy $D^q$ do środka, zachowując
$S^p\times D^q$. Iloraz $D^{p+1}/S^p$
jest sferą $S^{p+1}$, więc homologia względna
wynosi $\Z$ tylko w stopniu $p+1$.

Po odwróceniu śladu przyczepiamy do $M_e$
uchwyt indeksu $q$. Ten sam argument daje
$(W_e,M_e)\simeq(D^q,S^{q-1})$ na poziomie
homologii względnej i drugą formułę.
Ciąg dokładny pary, na przykład,
\[
 \cdots\longrightarrow H_{p+1}(M)
 \longrightarrow H_{p+1}(W_e)
 \longrightarrow\Z\xrightarrow{\partial}
 H_p(M)\longrightarrow H_p(W_e)
 \longrightarrow0,
\]
pokazuje konkretnie, że generator nowego
względnego $(p+1)$-dysku ma na brzegu klasę
rdzenia $[e(S^p\times\{0\})]$.
Analogiczny ciąg względem $M_e$ opisuje
sferę dualną. Nie wolno z samego wyrażenia
$H_{p+1}(W_e,M)=\Z$ wnioskować, że zawsze
ubywa generator w $H_p(M)$: jeśli klasa rdzenia
jest zerowa, nowy dysk może zamiast tego
tworzyć klasę w $H_{p+1}(W_e)$.
\end{proof}

\begin{stwierdzenie}[Zmiana charakterystyki Eulera]
\label{prop:euler-chirurgia-21}
Dla $M$ zamkniętej zachodzi
\begin{equation}\label{eq:euler-chirurgia-21}
 \chi(M_e)-\chi(M)
 =(-1)^{q-1}-(-1)^p
 =\begin{cases}
    0,&n\text{ nieparzyste},\\
    -2(-1)^p,&n\text{ parzyste}.
   \end{cases}
\end{equation}
\end{stwierdzenie}
\begin{proof}
Oznaczmy przez $C$ część $M$ pozostałą
po usunięciu wnętrza $S^p\times D^q$.
Z addytywności charakterystyki Eulera przy
sklejaniu wzdłuż podkompleksu wspólne wyrazy
$\chi(C)$ i $\chi(S^p\times S^{q-1})$ redukują się.
Pozostaje
\[
 \chi(M_e)-\chi(M)
 =\chi(D^{p+1}\times S^{q-1})
  -\chi(S^p\times D^q)
 =(1+(-1)^{q-1})-(1+(-1)^p).
\]
Ponieważ $q=n-p$, przy nieparzystym $n$
wykładniki $q-1$ i $p$ mają tę samą parzystość,
a przy parzystym przeciwną. To daje
ostatnią postać~\eqref{eq:euler-chirurgia-21}.
\end{proof}

\begin{twierdzenie}[Wpływ na $\pi_1$ przy kowymiarze co najmniej trzy]
\label{thm:pi1-chirurgia-21}
Jeśli $q=n-p\geq3$, to inkluzja $M_e\hookrightarrow W_e$
indukuje izomorfizm na $\pi_1$. Ponadto:
\begin{enumerate}[label=(\alph*)]
\item dla $p\geq2$ obie inkluzje brzegów indukują
izomorfizmy na $\pi_1$;
\item dla $p=1$ mamy
$\pi_1(M_e)\cong\pi_1(M)/\langle\!\langle[e(S^1)]\rangle\!\rangle$;
\item dla $p=0$ (obie kule leżą w spójnej $M$)
$\pi_1(M_e)\cong\pi_1(M)*\Z$.
\end{enumerate}
Przez $\langle\!\langle\gamma\rangle\!\rangle$
oznaczamy najmniejszą podgrupę normalną
zawierającą $\gamma$.
\end{twierdzenie}
\begin{proof}
Ślad ma typ homotopijny $M$ z jedną dołączoną
komórką wymiaru $p+1$: czynnik $D^q$ uchwytu
jest ściągalny, a mapa przyczepienia to rdzeń $S^p$.
Dla $p\geq2$ dołączenie komórki wymiaru
co najmniej $3$ nie zmienia grupy podstawowej
(twierdzenie van Kampena lub przybliżanie pętli
i homotopii przez komórki wymiaru $2$).
Dla $p=1$ dołączenie $2$-komórki dodaje relację
$[e(S^1)]=1$ wraz ze wszystkimi jej sprzężeniami,
co daje iloraz z punktu (b). Dla $p=0$
dołączamy odcinek z końcami w tej samej spójnej
składowej $M$. Po wyborze drogi łączącej końce
otrzymujemy dodatkową pętlę niezależną od
poprzednich, więc van Kampen daje iloczyn wolny
z $\Z$.

Od strony $M_e$ ślad ma jedną komórkę wymiaru
$q\geq3$, gdyż dualny uchwyt ma indeks $q$.
Takie dołączenie nie zmienia $\pi_1$.
Złożenie powyższych obliczeń dowodzi tez.
Bez warunku $q\geq3$ ostatnie zdanie może być
fałszywe: dualny $2$-uchwyt potrafi dodać relację.
\end{proof}

\begin{przyklad}[Sfery przechodzą w produkty sfer]
\label{prz:sfera-produkt-chirurgia-21}
Podzielmy $S^n$ jako brzeg
$D^{p+1}\times D^q$:
\[
 S^n=(S^p\times D^q)
 \cup_{S^p\times S^{q-1}}
 (D^{p+1}\times S^{q-1}).
\]
Wybierzmy standardową sferę $S^p\times\{0\}$
w pierwszym kawałku. Po jej chirurgii pozostaje
$D^{p+1}\times S^{q-1}$ i doklejamy jego
drugi egzemplarz wzdłuż brzegu.
Podwojenie $D^{p+1}$ jest $S^{p+1}$, zatem
\begin{equation}\label{eq:chirurgia-sfery-21}
 (S^n)_e\cong S^{p+1}\times S^{q-1}
 \qquad\text{dla standardowego obramowania.}
\end{equation}
Dla $p=0$ uzyskujemy $S^1\times S^{n-1}$:
operacja \emph{tworzy} generator $\pi_1$.
Dla $n=6$, $p=2$ otrzymujemy
$S^3\times S^3$, którego $H_3$ ma rangę $2$,
mimo że $H_3(S^6)=0$. Chirurgia nie zawsze
usuwa homologię; zależy od położenia rdzenia
i od stopnia operacji.
\end{przyklad}

\section{Odwzorowania normalne: jakie dane zachowujemy?}
\label{sec:normalne-chirurgia-21}

Dotąd zmienialiśmy rozmaitość bez wskazania celu.
W programie chirurgii mamy przestrzeń $X$ i chcemy
poprawić odwzorowanie $f:M\to X$. Najczytelniejsza
wersja wprowadzająca zakłada, że $M^n$ i $X^n$
są spójnymi zamkniętymi \emph{zorientowanymi}
rozmaitościami gładkimi. Wersje z brzegiem
i współczynnikami lokalnymi omówimy później.

\begin{definicja}[Stopień jeden i odwzorowanie normalne]
\index{stopien jeden i odwzorowanie normalne@Stopień jeden i odwzorowanie normalne}
\label{def:normal-map-21}
Odwzorowanie $f:M^n\to X^n$ ma \emph{stopień jeden},
jeśli $f_*[M]=[X]\in H_n(X;\Z)$.
Wybierzmy rzeczywistą wiązkę $\xi\to X$.
\emph{Odwzorowaniem normalnym stopnia jeden}
nazywamy dane $(f,b)$, gdzie $f$ ma stopień jeden,
a $b$ jest stabilnym izomorfizmem wiązek
\begin{equation}\label{eq:normal-map-21}
 b:\nu_M\oplus\varepsilon^a
   \xrightarrow{\ \cong\ }
   f^*\xi\oplus\varepsilon^b
\end{equation}
po dobraniu rang tak, aby obie strony były tej samej
rangi. Tutaj $\nu_M$ oznacza wiązkę normalną
zanurzenia $M$ w przestrzeni euklidesowej,
a $\varepsilon^j$ trywialną wiązkę rangi $j$.
Dwa zanurzenia dają ten sam stabilny typ $\nu_M$.
\end{definicja}

W dalszych rozdziałach dopuszczamy też cel
nieorientowalny. Wtedy wymagamy
$f^*w_1(TX)=w_1(TM)$, czyli zgodności
lokalnych systemów orientacji, i definiujemy
stopień jeden przez
$f_*[M]=[X]$ w
$H_n(X;\Z^{w_X})$; po lewej używamy
$[M]\in H_n(M;\Z^{w_M})$.
Symbol $\Z^{w_X}$ oznacza lokalny system,
w którym pętla odwracająca orientację
działa na $\Z$ przez mnożenie przez $-1$.
Gdy obie rozmaitości są zorientowane,
jest to dokładnie poprzednia definicja.
Ogólna wersja twierdzenia Walla i ciągu
chirurgii korzysta z tej konwencji.

Warunek o stopniu mówi, że $f$ przenosi klasę
fundamentalną bez wielokrotności. Dane $b$ pozwalają
porównywać kierunki normalne nad punktami $M$
z wiązką wybraną nad $X$ i kontrolować obramowania
podczas operacji. Nie każda mapa stopnia jeden
\emph{automatycznie} ma strukturę~\eqref{eq:normal-map-21}
dla z góry ustalonej $\xi$. Gdy $X$ jest
rozmaitością, naturalnym punktem odniesienia jest
$\xi=\nu_X$: identyczność $X\to X$ ma wtedy
oczywiste dane normalne.

\begin{lemat}[Rozszczepienie na homologię celu i jądro]
\label{lem:split-normalne-21}
Jeśli $f:M^n\to X^n$ ma stopień jeden, to
$f_*:H_i(M;\Z)\to H_i(X;\Z)$ jest surjekcją
w każdym stopniu. Istnieje naturalny prawy odwrotny
\[
 f_!=\operatorname{PD}_M\circ f^*
       \circ\operatorname{PD}_X^{-1}:H_i(X;\Z)\to H_i(M;\Z),
 \qquad f_*f_!=\id.
\]
W szczególności
$H_i(M;\Z)\cong f_!H_i(X;\Z)\oplus K_i(f)$,
gdzie $K_i(f)=\ker f_*$.
\end{lemat}
\begin{proof}
Niech $x\in H_i(X;\Z)$ i niech
$\alpha=\operatorname{PD}_X^{-1}(x)
\in H^{n-i}(X;\Z)$.
Z definicji dualności Poincarégo
$f_!(x)=[M]\frown f^*\alpha$.
Naturalność iloczynu z klasą fundamentalną daje
\[
 f_*f_!(x)
 =f_*([M]\frown f^*\alpha)
 =f_*[M]\frown\alpha
 =[X]\frown\alpha=x.
\]
Zatem $f_*$ ma prawy odwrotny. Każde
$z\in H_i(M)$ rozkłada się jawnie jako
\[
 z=f_!(f_*z)+\bigl(z-f_!(f_*z)\bigr),
\]
a drugi składnik leży w $K_i(f)$.
Część wspólna obu składników jest zerowa,
bo $f_*f_!=\id$. Nie używaliśmy tu jeszcze
izomorfizmu wiązek $b$.
\end{proof}

\begin{uwaga}[Czego nie dowodzi ten lemat?]
Znikanie wszystkich $K_i(f)$ dla \emph{zwykłych}
homologii nie wystarcza do stwierdzenia, że $f$
jest równoważnością homotopijną przy dowolnej
$\pi_1(X)$. Dla pełnego kryterium potrzebna
jest kontrola $\pi_1$ oraz homologii nakryć
uniwersalnych, czyli współczynników w
$\Z[\pi_1(X)]$; używa się wówczas odpowiedniej
wersji twierdzenia Whiteheada. To odróżnia
zadanie chirurgii od samego liczenia rang homologii.
\end{uwaga}

\begin{definicja}[Normalna kobordancja]
\index{normalna kobordancja@Normalna kobordancja}
\label{def:normal-bordism-21}
Dwa odwzorowania normalne $(f_0,b_0,\xi_0):M_0\to X$
i $(f_1,b_1,\xi_1):M_1\to X$ są \emph{normalnie
kobordantne}, jeśli istnieje zwarty gładki kobordyzm
$(W;M_0,M_1)$, odwzorowanie $F:W\to X$
o ograniczeniach $f_0,f_1$ oraz wiązka
nad $X\times I$ ze stabilnymi danymi normalnymi
nad $F$, której ograniczenia na końcach dają
$(\xi_0,b_0)$ i $(\xi_1,b_1)$ (z uwzględnieniem
normalnego kierunku kołnierza brzegu).
W szczególności wiązki $\xi_0$ i $\xi_1$
muszą być stabilnie izomorficzne; w klasyfikacji
nie ustalamy jednak z góry jednej wiązki $\xi$.
\end{definicja}

\begin{stwierdzenie}[Stopień nie zmienia się w kobordancji]
\label{prop:stopien-kobordancja-21}
Przy zgodnym wyborze orientacji $W$ i jego brzegów,
tak że
$\partial[W]=[M_1]-[M_0]$, to
$F_*[M_0]=F_*[M_1]\in H_n(X;\Z)$.
\end{stwierdzenie}
\begin{proof}
Na poziomie łańcuchów $F_*(\partial[W])$
jest brzegiem $F_*[W]$, więc przedstawia zero
w homologii $X$. Po podstawieniu
$\partial[W]=[M_1]-[M_0]$ otrzymujemy
$f_{1*}[M_1]-f_{0*}[M_0]=0$.
Zatem jeśli jedno ograniczenie ma stopień jeden,
to drugie także.
\end{proof}

\section{Kiedy chirurgia działa na odwzorowanie?}
\label{sec:normalny-krok-21}

Zwykła operacja~\eqref{eq:chirurgia-21} nie
wystarcza: po zmianie dziedziny musimy nadal
mieć odwzorowanie do $X$ i zgodne dane normalne.

\begin{twierdzenie}[Warunek przedłużenia odwzorowania]
\label{thm:przedluzenie-mapy-21}
Niech $e:S^p\times D^q\hookrightarrow M$ będzie
zanurzeniem obramowanym, $f:M\to X$ ciągłym
odwzorowaniem. Odwzorowanie $f$ przedłuża się
(po homotopii w otoczeniu $e$) do
$F:W_e\to X$ wtedy i tylko wtedy, gdy
\[
 f\circ e|_{S^p\times\{0\}}:S^p\to X
\]
jest nullhomotopijne. Aby przedłużyć także
\emph{dane normalne} $b$ z~\eqref{eq:normal-map-21},
wybrana nullhomotopia i obramowanie $e$
muszą być zgodne ze stabilną trywializacją
nad przyklejanym uchwytem. Dokładniej,
po ustaleniu trywializacji obu wiązek nad
kontraktybilnym uchwytem różnica między
identyfikacjami na części przyczepionej
$S^p\times D^q\simeq S^p$ daje stabilną mapę
$S^p\to O$. Dane normalne przedłużają się
wtedy i tylko wtedy, gdy ta mapa jest
nullhomotopijna.
\end{twierdzenie}
\begin{proof}
Jeżeli $F$ istnieje, ograniczmy je do rdzenia
uchwytu $D^{p+1}\times\{0\}\subset W_e$.
Na jego brzegu dostajemy $f\circ e|_{S^p\times\{0\}}$;
rdzeń jest więc żądaną nullhomotopią.

Odwrotnie, niech $H:D^{p+1}\to X$ przedłuża
to odwzorowanie rdzenia. W otoczeniu
$S^p\times D^q$ odwzorowanie $f\circ e$
jest homotopijne do mapy niezależnej od
współrzędnej $D^q$, przez skurczenie
$z\mapsto(1-t)z$. Własność przedłużania
homotopii dla kołnierza zanurzonej tuby
pozwala zastąpić $f$ mapą homotopijną,
która na obszarze przyczepienia zależy
tylko od $x\in S^p$. Na uchwycie kładziemy
$F(y,z)=H(y)$, $y\in D^{p+1}$, $z\in D^q$.
Formuły zgadzają się na $S^p\times D^q$,
więc sklejają się w mapę całego $W_e$.

Wiązka nad dyskiem $D^{p+1}\times D^q$ jest
trywialna; po stabilizacji wybieramy na niej
ramy. Obramowanie $e$ identyfikuje ramy
normalne na obszarze przyczepienia z ramami
pochodzącymi od $M$. Ich porównanie za pomocą
$b$ jest mapą $S^p\times D^q\to O(r)$
dla dostatecznie dużego $r$. Ponieważ $D^q$
jest ściągalny, mapa ta ma klasę w $\pi_p(O)$.
Przedłużenie na $D^{p+1}\times D^q$ istnieje
dokładnie wtedy, gdy jej ograniczenie do $S^p$
jest nullhomotopijne; wystarczalność daje
rozciągnięcie nullhomotopii radialnie po
$D^{p+1}$, a konieczność jej ograniczenie
z dowolnego przedłużenia.
Inny wybór ram nad dyskiem zmienia porównanie
o mapę przedłużalną, zatem nie zmienia warunku.
Sam fakt, że $H$ istnieje, nie zeruje tej
klasy: warunek normalny jest dodatkowy.
Wtedy wraz z $F$ otrzymujemy normalną
kobordancję do odwzorowania $M_e\to X$.
\end{proof}

\begin{figure}[htbp]
\centering
\begin{tikzpicture}[>=Latex,font=\small]
 \node[draw=manifoldblue!75!black,fill=manifoldblue!9,
    rounded corners=3pt,minimum width=1.9cm,minimum height=.65cm]
    (S) at (-3.35,.92) {$S^p$};
 \node[draw=manifoldgreen!70!black,fill=manifoldgreen!9,
    rounded corners=3pt,minimum width=1.9cm,minimum height=.65cm]
    (M) at (-.25,.92) {$M$};
 \node[draw=manifoldgold!75!black,fill=manifoldgold!11,
    rounded corners=3pt,minimum width=1.9cm,minimum height=.65cm]
    (D) at (-3.35,-.95) {$D^{p+1}$};
 \node[draw=manifoldred!75!black,fill=manifoldred!8,
    rounded corners=3pt,minimum width=1.9cm,minimum height=.65cm]
    (X) at (-.25,-.95) {$X$};
 \draw[->,thick] (S)--node[above] {$e|_{S^p}$} (M);
 \draw[->,thick] (M)--node[right] {$f$} (X);
 \draw[->,thick] (S)--node[left] {brzeg} (D);
 \draw[->,thick,manifoldred!75!black]
    (D)--node[below] {$H$} (X);
 \node[align=left,text width=3.7cm] at (3.3,0)
  {Dysk $H$ wypełnia obraz sfery w $X$.\\
   Zgodne obramowanie pozwala\\
   dołączyć cały uchwyt nad $H$.};
\end{tikzpicture}
\caption{Warunek na odwzorowanie w kroku chirurgii.
Komutatywność diagramu oznacza
$H|_{S^p}=f\circ e|_{S^p\times\{0\}}$.
Dodatkowe dane normalne kontrolują obramowanie.}
\label{fig:normalna-chirurgia-21}
\end{figure}

\begin{przyklad}[Chirurgia nad sferą docelową]
\label{prz:normalna-sfera-21}
Dla standardowej sfery $S^p\subset S^n$, $p<n$,
identyczność $S^n\to S^n$ wysyła ten rdzeń
na sferę nullhomotopijną, ponieważ
$\pi_p(S^n)=0$. Przy standardowym obramowaniu
i zgodnych stabilnych danych normalnych
krok chirurgii daje mapę stopnia jeden
\[
 f':S^{p+1}\times S^{q-1}\longrightarrow S^n
\]
normalnie kobordantną z identycznością.
Możemy też zbudować $f'$ bezpośrednio:
wybieramy mały $n$-dysk w produkcie,
przenosimy go na $S^n$ z zachowaniem
orientacji i zgniatamy do punktu dopełnienie
wnętrza dysku. Otrzymana mapa ma stopień jeden.
Produkt sfer jest stabilnie równoległizowalny
(po dodaniu dwóch prostych trywialnych),
a $S^n$ po dodaniu jednej; dlatego istnieją
zgodne stabilne dane normalne.
Nie twierdzimy, że $f'$ jest równoważnością:
dla $1\leq p\leq n-2$ produkt ma zwykle
niezerową homologię w stopniach pośrednich.
\end{przyklad}

\section{Jądra chirurgii i próg wymiaru środkowego}
\label{sec:jadra-chirurgii-21}

Dla odwzorowania normalnego stopnia jeden
piszemy $K_i(f)=\ker(H_i(M;\Z)\to H_i(X;\Z))$.
Intuicyjnie są to klasy, które dziedzina ma
\emph{nadmiarowo} w porównaniu z celem.
Nie każdą z nich można usunąć: trzeba ją
zrealizować przez zanurzoną sferę z obramowaniem,
po której obraz w $X$ da się wypełnić dyskiem.
Obramowanie musi ponadto pasować do danych
normalnych. W wyższych wymiarach pozycja ogólna
i trik Whitneya pozwalają realizować wiele
takich operacji poniżej wymiaru środkowego;
przy wymiarze środkowym sfery spotykają się
i pojawia się dodatkowy rachunek przecięć.

\subsection*{Co dokładnie usuwa trik Whitneya?}
\index{trik Whitneya@Trik Whitneya}
Niech $A^p$ i $B^q$ będą zorientowanymi,
poprzecznie zanurzonymi podrozmaitościami
w zorientowanej $N^n$, gdzie $p+q=n$.
Zakładamy $p,q\geq1$.
Ich przecięcia są izolowanymi punktami ze znakami.
Wybierzmy dwa punkty $x_-,x_+\in A\cap B$
o przeciwnych znakach oraz zanurzone łuki
$a\subset A$ od $x_-$ do $x_+$
i $b\subset B$ wracający do $x_-$.
Ich wnętrza omijają drugi płat i wszystkie
inne punkty przecięcia, a $a\cap b=\{x_-,x_+\}$.
Tworzą \emph{okrąg Whitneya} $a\cup b$.
Samo równanie $(-1)+(+1)=0$ nie daje jeszcze
izotopii: potrzebny jest dysk wypełniający ten
okrąg w odpowiednim położeniu.

\begin{stwierdzenie}[Lokalny ruch Whitneya]
\label{prop:lokalny-Whitney-21}
Załóżmy, że $a\cup b$ ogranicza gładko zanurzony
dysk $D_W\subset N$ o wnętrzu rozłącznym z
$A\cup B$. Załóżmy także, że dysk ma
\emph{obramowanie Whitneya}: normalne kierunki
indukowane wzdłuż $a$ i $b$ przez oba płaty
przedłużają się bez skręcenia nad $D_W$.
Wówczas lokalna izotopia jednego z płatów,
wsparta w otoczeniu $D_W$, usuwa $x_-,x_+$
bez tworzenia nowych przecięć $A\cap B$.
\end{stwierdzenie}
\begin{proof}
Obramowanie daje produktowe otoczenie dysku.
W nim łuk $a$ leży na pierwszym płacie,
a łuk $b$ na drugim; przy dwóch końcach
oba płaty przecinają się poprzecznie.
Przesuwamy wąski pasek pierwszego płata
wzdłuż $D_W$ na drugą stronę drugiego płata.
W przekroju prostopadłym do paska jest to
przejście przez dysk ograniczony przez $a\cup b$.
Przeciwne znaki końców gwarantują, że oba
lokalne przecięcia znikają jednocześnie.
Produktowe otoczenie i rozłączność wnętrza
$D_W$ z płatami pozwalają wykonać ruch bez
nowych przecięć. Poza tym otoczeniem płaty
pozostają na miejscu; przedłużenie izotopii
z Sekcji~\ref{sec:izotopie-otoczen-15}
daje izotopię otoczenia w $N$.
\end{proof}

\begin{figure}[htbp]
\centering
\begin{tikzpicture}[>=Latex,font=\small,line cap=round,
 line join=round]
 \fill[manifoldgold!17]
   (-1.05,0) .. controls (-.52,-.91) and (.52,-.91)
       .. (1.05,0) -- cycle;
 \draw[manifoldgold!75!black,densely dashed,thick]
   (-1.05,0) .. controls (-.52,-.91) and (.52,-.91)
       .. (1.05,0);
 \draw[manifoldblue!85!black,very thick] (-2.55,0)--(2.55,0);
 \draw[manifoldred!82!black,very thick]
   (-2.35,.65) .. controls (-1.70,.72) and (-1.28,.37)
       .. (-1.05,0) .. controls (-.52,-.91) and (.52,-.91)
       .. (1.05,0) .. controls (1.28,.37) and (1.70,.72)
       .. (2.35,.65);
 \fill[manifoldblue!90!black]
   (-1.05,0) circle (2pt) (1.05,0) circle (2pt);
 \node[above] at (-1.05,.10) {$x_-$};
 \node[above] at (1.05,.10) {$x_+$};
 \node[manifoldblue!85!black] at (-2.38,-.30) {$A$};
 \node[manifoldred!82!black] at (-2.20,.96) {$B$};
 \node[manifoldgold!65!black] at (0,-.33) {$D_W$};
 \draw[->,thick,manifoldblue!70!black] (2.78,.05)--(3.74,.05);
 \begin{scope}[shift={(6.25,0)}]
  \draw[manifoldblue!85!black,very thick] (-2.15,0)--(2.15,0);
  \draw[manifoldred!82!black,very thick]
    (-2.05,.65) .. controls (-.80,.76) and (.80,.76)
      .. (2.05,.65);
  \node[manifoldblue!85!black] at (-1.98,-.30) {$A$};
  \node[manifoldred!82!black] at (-1.90,.96) {$B$};
 \end{scope}
\end{tikzpicture}
\caption{Przekrój schematyczny ruchu Whitneya.
Po lewej łuki na $A$ i $B$ ograniczają jasny dysk,
a $x_-$ i $x_+$ mają przeciwne znaki. Po ruchu
przecięć nie ma. Rysunek pokazuje tylko przekrój;
w wyższych wymiarach niezbędne jest obramowanie
całego dysku.}
\label{fig:Whitney-chirurgia-21}
\end{figure}

Istnienie dysku z założeń stwierdzenia jest
\emph{osobną kwestią}. Okrąg $a\cup b$ musi być
nullhomotopijny; nawet wtedy dysk początkowo
może przecinać płaty lub mieć złe obramowanie.
Jeśli $\pi_1(N)\ne1$, dwa punkty o przeciwnych
znakach mogą mieć różne etykiety dróg $g,h$,
więc dają $g-h\ne0$ w $\Z[\pi_1(N)]$.
W wymiarze środkowym $n=2k$, $k\geq3$,
pełne kryterium zamiany obramowanej immersji
$S^k\looparrowright N^{2k}$ na zanurzenie
wyraża kwadratowe samoprzecięcie $\mu$;
to \href{https://webhomes.maths.ed.ac.uk/~v1ranick/books/scm.pdf}
{twierdzenie 5.2 Walla}, a nie wniosek
z samego rysunku. Dla $k=2$ kryterium to
może zawodzić. Tak właśnie powstaje granica
między chirurgią wysokowymiarową a wymiarem $4$.

\begin{stwierdzenie}[Forma na jądrze w wymiarze parzystym]
\label{prop:forma-jadra-21}
Niech $n=2k$, a $f:M^{2k}\to X^{2k}$ ma stopień
jeden. Forma przecięć na środkowej homologii
$H_k(M;\Z)/\mathrm{tors}$ rozkłada się ortogonalnie
na obraz $f_!$ oraz $K_k(f)/\mathrm{tors}$.
Forma na jądrze jest niezdegenerowana
(nad $\Z$: unimodularna na częściach wolnych).
Dla parzystego $k$ jest symetryczna, a dla
nieparzystego antysymetryczna.
\end{stwierdzenie}
\begin{proof}
Na częściach wolnych dualność Poincarégo
identyfikuje formę przecięć z iloczynem kubkowym:
$\lambda_M(x,y)=\langle\operatorname{PD}_M^{-1}x
\smile\operatorname{PD}_M^{-1}y,[M]\rangle$.
Dla $a,b\in H_k(X)$ z naturalności i warunku
$f_*[M]=[X]$ dostajemy
\[
 \lambda_M(f_!a,f_!b)
 =\langle f^*\operatorname{PD}_X^{-1}a
 \smile f^*\operatorname{PD}_X^{-1}b,[M]\rangle
 =\lambda_X(a,b).
\]
Dla $z\in K_k(f)$ oraz $a\in H_k(X)$,
naturalność iloczynu z klasą fundamentalną
i definicja $f_!$ dają
$\lambda_M(z,f_!a)=\lambda_X(f_*z,a)=0$.
Lemat~\ref{lem:split-normalne-21} daje rozkład
$H_k(M)=f_!H_k(X)\oplus K_k(f)$, więc jest
on ortogonalny. Na częściach wolnych forma
całej zamkniętej zorientowanej rozmaitości
jest unimodularna; podobnie forma na $X$.
Jeżeli w bazach dostosowanych do rozkładu
macierz jest blokowo diagonalna, a wyznaczniki
całości i bloku $X$ wynoszą $\pm1$, to także
wyznacznik bloku jądra wynosi $\pm1$.
Znak zamiany argumentów wynika z
$\alpha\smile\beta=(-1)^{k^2}\beta\smile\alpha$.
\end{proof}

\begin{przyklad}[Dwie sfery i ich jedno przecięcie]
\label{prz:forma-produkt-21}
Dla $M=S^k\times S^k$, $X=S^{2k}$, $k\geq3$,
weźmy mapę stopnia jeden z
Przykładu~\ref{prz:normalna-sfera-21}
(w tym kierunku wystarczy mapa zgniatająca).
Wtedy $K_k(f)=H_k(M)\cong\Z^2$,
generowane przez
$A=S^k\times\{*\}$ i $B=\{*\}\times S^k$.
Obie sfery mają trywialne wiązki normalne.
Możemy je przesunąć od samych siebie, zatem
$A\cdot A=B\cdot B=0$, ale $A$ i $B$
przecinają się w jednym punkcie.
W bazie $(A,B)$ otrzymujemy
\[
 \lambda=\begin{pmatrix}0&1\\(-1)^k&0\end{pmatrix}.
\]
Gdy $k$ jest parzyste, forma jest symetryczna;
gdy nieparzyste, antysymetryczna.
Nie wystarczy osobno policzyć, że obie sfery
reprezentują klasy jądra: ich wzajemne
przecięcie kontroluje możliwość kolejnych
operacji. W bardziej ogólnych sytuacjach
potrzebna jest ponadto informacja o
samoprzecięciach i obramowaniach, prowadząca
do form kwadratowych teorii chirurgii.
\end{przyklad}

\begin{uwaga}[Dwie różne przeszkody]
Grupy $L_n(\Z[\pi])$ opisują przeszkody
\emph{wykonania} dalszej chirurgii na mapie
normalnej, szczególnie w wymiarze środkowym.
Torsja Whiteheada w
$\operatorname{Wh}(\pi)$, poznana w
Rozdziale~\ref{ch:s-kobordyzm}, kontroluje
\emph{prostotę} otrzymanej równoważności
i iloczynowość $h$-kobordyzmu.
Są to różne etapy argumentu. Także późniejszy
ciąg dokładny chirurgii będzie wymagał
rozróżnienia zwykłych i prostych
równoważności homotopijnych.
\end{uwaga}

\section{Co pozostaje do zrobienia w programie chirurgii}
\label{sec:plan-chirurgia-21}

Dotychczasowe kroki można streścić bez
ukrywania danych potrzebnych dalej:
\begin{enumerate}[label=(\arabic*)]
\item Zaczynamy od odwzorowania normalnego
stopnia jeden $(f,b):M^n\to X^n$.
\item Wybieramy obramowane sfery reprezentujące
klasy jądra i wykonujemy chirurgię nad $X$,
pozostając w tej samej klasie normalnej kobordancji.
\item Sprawdzamy grupę podstawową i homologiczny
kompleks jądra nad $\Z[\pi_1(X)]$.
W wymiarze środkowym samo znikanie zwykłych
homologii już nie wystarcza: decydują formy
przecięć i ich kwadratowe uzupełnienia.
\item Po usunięciu tej przeszkody otrzymujemy
równoważność homotopijną; do klasyfikacji
rozmaitości trzeba jeszcze zrozumieć
jej torsję oraz klasę odpowiedniego
kobordyzmu.
\end{enumerate}

Następny rozdział zdefiniuje grupy chirurgiczne
$L_n(\Z[\pi])$ w odpowiednich wersjach,
pokaże proste przykłady form i sformułuje
twierdzenie o przeszkodzie chirurgicznej.
Później zbierzemy dane w ciągu dokładnym
chirurgii i zastosujemy je do klasyfikacji
rozmaitości.

\par\smallskip
{\footnotesize\noindent\emph{Dalsza lektura:}
\href{https://link.springer.com/chapter/10.1007/978-3-031-56334-8_4}
{W.~Lück i T.~Macko, \emph{The Surgery Step and $\xi$-Bordism}},
\href{https://webhomes.maths.ed.ac.uk/~v1ranick/books/scm.pdf}
{C.~T.~C. Wall, \emph{Surgery on Compact Manifolds}, §1--3}.\par}
